% ============================================================
%  Poul Martin Møller,
%  "Ontology, or the System of the Categories"
%  ("Ontologien eller Kategoriernes System")
%
%  English translation of the fragment (Brudstykke) published in
%  Efterladte Skrifter af Poul M. Møller, bd. 3 (Tredie Bind),
%  ed. F.C. Olsen. Copenhagen: C.A. Reitzel.
%  Printed pp. 189--212.
%
%  Translated from transcription.tex (the Danish source of truth).
%  Conventions: Danish „…" -> ``…''; \emph{} preserved; Latin ->
%  \textit{}; Greek copied verbatim; the old "id est" mark ɔ: ->
%  "i.e."; section-break rules reproduced. Printed page numbers ->
%  \opage{N}.
% ============================================================

\documentclass[12pt]{article}
\usepackage[utf8]{inputenc}
\usepackage[T1]{fontenc}
\usepackage{libertinus}
\usepackage{libertinust1math}
\usepackage{textalpha}
\usepackage[english]{babel}
\usepackage[protrusion=true,expansion=true]{microtype}
\usepackage[margin=1.4in]{geometry}
\usepackage{setspace}
\usepackage{fancyhdr}
\usepackage{marginnote}
\usepackage{enumitem}
\usepackage[colorlinks=true,linkcolor=black,urlcolor=black,
  pdftitle={Ontology, or the System of the Categories},
  pdfauthor={Poul Martin Møller}]{hyperref}
\setlength{\headheight}{15pt}
\pagestyle{fancy}
\fancyhf{}
\fancyhead[LE,RO]{\thepage}
\fancyhead[RE]{\textit{Ontology, or the System of the Categories}}
\fancyhead[LO]{\textit{Poul Martin Møller}}
\setlength{\parindent}{1.5em}
\setlength{\parskip}{0pt}
\onehalfspacing

\newcommand{\opage}[1]{\marginnote{\footnotesize\textit{[#1]}}}

\title{\textbf{Ontology}\\[0.3em]
  \large or\\[0.3em]
  \textbf{The System of the Categories}\footnote{This is how the author
  designates the science he lectured on, in the announcement of his lectures on
  it in the winter half-year of 1837.}}
\author{Poul Martin Møller}
\date{Efterladte Skrifter, vol.\ 3}

\begin{document}
\maketitle

\begin{center}\textsc{Introduction.}\end{center}

\opage{189}The introduction to metaphysics, in the proper sense of the word, is a
peculiar philosophical discipline, one that develops the various stages through
which human consciousness runs from its immediate shape up to the level of
cultivation at which metaphysics takes its beginning. Such a philosophical
propaedeutic Hegel gave in his Phenomenology of Spirit; but that work, through
the manifoldness of its content, far oversteps the boundaries marked out by that
definition, and moreover presupposes readers who have an acquaintance with the
history of philosophy. Other philosophers too have, under various names, given
presentations of the philosophical propaedeutic in that sense. The most
distinguished are due to J.\ E.\ von Berger, C.\ F.\ Krause, and J.\ H.\ Fichte.
In a certain sense, then, metaphysics cannot be called the beginning of
philosophy, since the propaedeutic must, in a suitably arranged exposition of the
system of philosophy, be sent on ahead. This \opage{190}science is meant, namely,
to lead the reader to the boundaries of metaphysics and so to orient him that he
clearly recognizes into which region he is now to be conducted. It is to teach
him to understand what is meant when the thinking he is now called upon to follow
is termed immanent or a priori thinking. With respect to the natural order of the
didactic exposition, the propaedeutic is thus to be regarded as the first
discipline of the system of philosophy; on another consideration, metaphysics
must be called the beginning of philosophy. For it proceeds from the most
universal determinations of thought valid for all reality, determinations that
must be presupposed in the more concrete concepts gone through in the
propaedeutic. In concept, therefore, metaphysics is the truer beginning, to which
that preparatory science goes back.

Metaphysics, then, can not only, like every philosophical discipline, be
expounded on the strength of its relative self-sufficiency in such a way that,
taken by itself, it becomes intelligible; it can, with greater right than any
other part of philosophy, in its isolated shape make a claim to intelligibility.
For its determinations are repeated, in more concrete shapes, in all the parts of
philosophy, while it itself contains the simplest elements of all cognition. When
the propaedeutic, by leading up to the region in which metaphysical thinking
moves, has performed its service, the propaedeutic may be lost from view like a
ladder that has fulfilled its purpose. The new thinking to which one now turns is
namely no longer conditioned by that introduction, but comes forward as wholly
presuppositionless and immediate. As a consequence, the thoughtful reader who
follows the exposition of metaphysics with his attention (even if he does not, at
its very beginning, at once notice in what sphere he finds himself) will, in the
course of its advance, come to understand the light that the various propositions,
by their mutual connection, cast upon one another.

\opage{191}On the few pages to which I have here given the name of Introduction,
it is naturally not my intention to follow the train of thought demanded by a
rigorous, scientific propaedeutic. I merely send ahead a couple of aphoristic
remarks, by which I believe I may awaken in one reader or another a preliminary
representation of the content of metaphysics and of its significance within the
whole system of science. Whether these introductory considerations meet with
points of contact in the reader's circle of consciousness, or not, is entirely
accidental. Of them holds what holds of aphorisms in general: they are like hats
made without measurement; it is a stroke of luck whether any of them fits a given
individual's head, or not.

If, for one acquainted with the critical philosophy, one were to set out a
preliminary concept of the principal part of metaphysics that is to be expounded
here, one could give him the information that it aimed at a systematic deduction
of the categories; for, with the addition of a few closer determinations, that
definition would give him a tolerable representation of the content of the science
(of the writing). In the critical philosophy one strove, namely, to give an
inventory of the concepts that lie at the ground of thinking and that are not
given through sensory sensation. These concepts Kant sought to bring to
consciousness by reflection on the various forms of judgments. For this
investigation he laid at the ground the division of judgments familiar from the
usual logics (namely the one that takes account of their quality, quantity,
modality, and relation). It was taken in his philosophy as a matter of fact that
those concepts are presupposed in all thinking that posits objective reality. But
he makes no attempt to show how it lies in the nature of the case that thinking
comes forward precisely in these forms. Nor does he give any particular account
of these concepts' mutual relation to one another; although the question would
naturally have to \opage{192}be raised whether these concepts are merely
coordinate with one another without further connection, or whether they stand to
one another as universal and particular. --- Kant himself gives occasion for this
question, in that he, in passing, loosely throws out the remark that the third
category in every class arises through a combination of the first and the second.
But this loosely thrown-out thought, on which Kant himself lays no weight, prompts
one strongly to ask further what relation the various classes of categories then
have to one another. Such riddles are voiced, more or less definitely, all through
the development of the critical philosophy, where they were summed up in the
definite demand for a deduction of the categories.

But in taking this refuge in the idiom of the critical philosophy, it is of
particular importance that one strike out of one's head every representation of
the subjectivity of the categories. Kant assumed, namely, that these concepts,
which lie at the ground of thinking, are merely forms that are necessary for the
human being if he is to be able to ascribe objective reality to his sensory
representations, without these therefore holding, as true determinations, for the
thing in itself. This purely arbitrary assertion is recognized as invalid in a
correct presentation of ontology.

The determinations that, in the thinking of existence, carry with them complete
conviction hold necessarily for the thing in itself. If one does not accept this
proposition as correct, one is entangled in a doubt that, consistently pursued,
cancels itself. One cannot, indeed, refute the skeptic who, at every step of a
scientific train of thought, stubbornly holds fast to the remark that human
thinking is bound to subjective laws, in accordance with which it forms for
itself only crooked concepts of the nature and essence of things; for against
every conceivable proof he need only utter his formula, ready once and for all:
``Thus I am compelled, according to my subjectivity, to think \opage{193}to
myself the essence of things; but who can vouch to me that I am not, without
knowing it, bound to the fate of being led to merely errant thoughts that
nevertheless, for me, carry complete conviction.'' --- Science can indeed utter,
for that absolute skeptic, a worldview according to which the conviction of the
reality of thinking is strengthened by a concept of the significance of thinking
within existence and of its springing from the object of religion, the
all-comprehending Essence; but the skeptic can, once the concept has been uttered,
come forward anew with his simple formula: ``Thus we must think to ourselves the
fundamental relations of existence, but perhaps in themselves they are not such as
we must think them.'' --- We are then obliged to answer that skeptic, who believes
that with a single magic formula he can annihilate all the fruits of scientific
thinking: Let it, then, be possible that the use of reason leads to mere
delusion --- a possibility that, taken by itself, can as little be disputed as
many another, e.g.\ the one that the moon is a glowing stone --- yet we
philosophize under the presupposition that what, for a rational thinking, carries
conviction holds for the thing in itself. We assume, then, once and for all, that
what rational beings must necessarily think to be, necessarily is, just as it must
be thought; and we desire no surer science than the one that, for us, carries
complete conviction. --- After these restrictions and corrections, we may
provisionally let that definition stand, that ontology is a deduction of the
system of the categories.

For a natural philosopher of the Schellingian school, ontology is best explained
as the elucidation (development) of the concept of the Absolute. For in
Schelling's philosophy this concept appears without any introduction, and the few
determinations by which it is elucidated are wholly unintelligible to one who has
no acquaintance with the philosophical movements that preceded \opage{194}natural
philosophy. Only through this historical knowledge can the concepts
\emph{intellectual intuition, the Absolute, subjectivity, objectivity, etc.}, of
which Schelling makes an immediate use, have any definite meaning for the reader.
This lack has also often enough been marked with the apt expression that
metaphysics or ontology was wanting among the natural philosophers.

The various definitions of the Absolute, found here and there in Schelling's many
writings, would naturally have to be brought together by the attentive reader. The
Absolute is now designated as the identity of the ideal and the real, now as the
identity of unity and opposition, now as the bond between unity and manifoldness,
now as the identity of essence and form, or of the affirming and the affirmed, and
so forth. But there is no clear consciousness at all of the meaning of these
various fundamental concepts and of their relation to one another. And yet it is
sometimes clear from the context that it is the nature of the matter itself that
brings Schelling, in new regions of his system, suddenly to alter the terminology
he had hitherto used. To this indeterminacy and instability in the designation of
the concepts there is added, moreover, the awkward circumstance that the names of
natural phenomena are frequently used in place of pure determinations of thought,
as when the words \emph{Light} and \emph{Gravity} are not taken in their usual
meaning but are used to designate fundamental concepts within the system.

Since the age no longer found itself satisfied with the half-poetic form in which
Schelling had presented his intuition of the nature and essence of the Absolute
(especially after a swarm of imitators had brought the defects of the presentation
to clear consciousness), G.\ W.\ Hegel attacked and carried out the great work of
developing all the fundamental concepts of philosophy in systematic connection. By
his labors the idea of ontology, in the form that the present shape of philosophy
requires, has been reborn in the philosophical \opage{195}world. The consequence
of this has also been that the several distinguished philosophers who, in our
time, have given presentations of this reawakened science have, more or less, made
use of his preliminary labors and formed themselves after the great model he set
up in his Wissenschaft der Logik.

The aim of ontology, then, becomes that of bringing to scientific clarity all the
concepts of which philosophy must make use in the complete definition of the
Absolute. The Absolute is therein to be considered according to the universal
determinations of its content, without regard to the special forms under which it
comes forward in time and space. All the concepts that, according to their deeper
meaning, can be thought as essential predicates of the Absolute are to be grasped
according to this meaning and to yield a contribution toward a construction of
cognition in the light of the Absolute.

We have said above that ontology is a self-contained circle of thought which, by
a clear presentation, can, like every other philosophical discipline, be made
intelligible taken by itself, although greater clarity is shed upon it when it is
seen in its organic connection with the whole system of philosophy. But though
ontology has such a relative self-sufficiency that it can be understood of itself,
in that its various parts cast the necessary light upon one another, it is not
thereby said that a reader who is devoid of all philosophical culture at once
grasps its very first propositions. If such a reader at once, at the beginning of
the exposition, finds himself oriented and notices in which region of thought he
finds himself, this must be regarded rather as a stroke of luck than as an effect
that can be counted on and that lies in the nature of the matter.

But it is not likely that this writing will be read by very many who wholly lack
acquaintance with the elements of philosophy. On the contrary, I hold it probable
that most of those who will lend it any attention have read or
stu\opage{196}died Professor Sibbern's most important works. For such readers a
few more considerations are sent ahead, in which but one concept is laid at the
ground as wholly familiar --- I mean the concept: a priori concepts, or
categories. It is here taken as a settled matter that at the ground of all
cognition there lies a portion of concepts that do not present any content given
through sensory sensation. Every kind of thinking, in the broad sense of the word,
has at every step of its development such categories among its presuppositions,
and applies them, even without consciousness of it, in every one of its movements.
We have said that these concepts lie at the ground of all cognition. This
proposition has quite literal truth and complete universal validity; for all
cognition comes to actuality as thinking, and thinking consists in an application
of a priori presuppositions that spring solely from thinking itself and lead its
activity forward. By these categories, determinations are cognized in actuality
which are there only for thinking, and which yet, for the thinking subject, carry
the same conviction of objective validity as the determinations that, in
consequence of affections of the nervous system, are ascribed to things. Thinking
pervades cognition at every one of its stages of development: not even a single
act of outward sensing can take place except in such a way that the subject
distinguishes something existing from itself and posits it as something for
itself, distinct from another existing thing; but here at once several categories
come into application, namely the concepts: existence, being-for-itself, etc.

There is, in the development of consciousness, a period in which the human being
thinks according to his categories with complete assurance, without misgiving or
doubt about their proper meaning. In this state of innocence of thinking, these
concepts constitute the firm supports about which the development of cognition
turns, but which are themselves exempt from the unrest of the movement. Even after
the human being has, by \opage{197}reflection on the presuppositions of thinking
and by imitation of foreign speech that goes beyond his understanding, lost that
sure tact for their proper use, a good portion of the most abstract categories
still retain their full meaning for him. Nor was it conceivable that all the
categories in a human being's consciousness should at once be brought into
confusion and dissolution without derangement being the inevitable consequence
thereof.

So long as thinking retains that merely natural soundness, consciousness can
naturally answer more or less to its idea. The clarity with which the a priori
concepts come to development in the individual at this standpoint depends first of
all on the richness and definiteness that prevails in the vocabulary of the mother
tongue, and next on the shape under which language first comes to his ear in his
circle of intercourse. One who is awakened to self-consciousness by a language so
poor in determinations of thought as the Greenlandic is set, by an unavoidable
fate, on a low level of intellectual culture, above which his individual natural
gifts are not able to raise him very far.

We are speaking here of the categories as they first show themselves, when they
are merely immediately held fast in definite forms within language. They are then
best regarded as a tradition, as the common property of a tribe of people, as its
intellectual patrimony and one of the most characteristic expressions of the whole
people's power of thought. The categories come, indeed, first to the individual as
a tradition within language, but are not therefore taken up by it as something
foreign. When this tradition is appropriated by a well-endowed individual, it
finds therein, on the contrary, a form for its own thinking, which it uses as
surely as if it had itself produced it. It thinks independently in the forms that
belong to the tribe, forms that are no more the one individual's property than the
other's, because they are a form of the universal reason, which has actuality only
in that it repeats itself in the thinking of the particular individuals.
\opage{198}The individual thinks for himself when he thinks in his tribe's forms
of thought; for his rational life consists only in connection with the rational
life of a society, whose great stream flows also through his individual
consciousness.

In the application of the categories, and in reflection on the categories in its
consciousness that have not lost their immediate meaning, the thinking subject
feels itself in its own element, more than in items of knowledge about the sensory
determinations of things: of this an immediate consciousness must be able to
convince every thinking subject. When it receives the knowledge that a thing is
red or sour, this knowledge is outward for it in a different sense than when it
recognizes that one phenomenon is cause and another effect. Thinking feels itself
otherwise disposed toward itself when it reflects on the concepts unity,
plurality, ground, consequence, etc., than when it considers the sensory
predicates ascribed to existing things. The first-named concepts are, namely,
thinking's own products; it recognizes in them the firm traces of its own
activity; it feels itself here at home in a different way than when it has to do
with representations and concepts that present merely the content of sensory
affections.

The categories can come forward either in a vulgar way or with a higher meaning.
In the first place they can lie at the ground of thinking in the period of merely
natural consciousness. They are then used, as the turn comes to them, each by
itself, according to a sound tact of the thinking subject, while it grasps its
circle of experience in a sensory and understanding way, or utters its religious
knowledge in a popular form. It is unable to form for itself representations of
the objects that, in the surrounding world, present themselves to its senses,
without its intuition constantly entering into union with a thinking that proceeds
from the categories as its firm presuppositions. Here \opage{199}the a priori
concepts show themselves precisely as wholly clear and intelligible to thinking;
for although the concepts existence, cause, effect, etc.\ here continually come
into application, no need is felt at this standpoint for their explanation: in this
period one never asks what existence, cause, effect mean; it is here understood of
itself; for one has to do, in these concepts, with the means of one's own
understanding. One thinks merely of understanding things, not of understanding
one's own understanding. Thus one has often, here in Denmark, even quite recently,
carried on a heated dispute about the proper boundary between rote learning and
understanding in popular religious instruction, without the disputing parties'
having expressed the least doubt that they understood one another's concept of
understanding. On the contrary, they had with assurance laid this concept at the
ground as a firm, common point of departure for their inquiries. Even in the
period in which the system of the categories is disturbed by reflections and is
in part held fast by isolated definitions, they can nevertheless still stand in the
service of vulgar thinking. This they do, namely, so long as they are used, as
isolated, firm points of departure, in the grasping of the phenomena of the given
actuality.

A deeper meaning a category acquires when it is grasped as a determination for all
actuality. Such a universal meaning one finds that now one, now another category
has won in the history of science. Sometimes, namely, one individual or another has
suddenly seen the validity of such a concept for everything that is, when it is
thought in its totality. It has, by such a single determination, embraced the whole
of actuality in its thinking, and held that it had thereby fully grasped it in its
truth. Thus several philosophers have with enthusiasm defended the proposition that
All was One, and have lived in the conviction that in this concept of unity they
had found the key to the whole riddle of existence. But such a \opage{200}one-sided
doctrine has usually called forth an equally forceful presentation of the
proposition, for an outward consideration wholly opposed to it, that what in truth
is is an infinite manifoldness. (We choose here, namely, an example that, as is
well known, was already brought forward by the Xenophontic Socrates, when he with
few words emphasized the great contradiction that obtained among the principles of
the philosophers.) --- But when a thinker holds that in a single category he has
discovered the true concept of actuality, there lies in his opinion both truth and
error. It is truth that a single category can be so thought that one has in it a
moment of the concept of actuality, which can be laid at the ground of a worldview.
It is incorrect, on the other hand, that he denies to other categories the same
universal meaning. The expressly uttered category designates only one of the sides
of a priori cognition; but he who, in such a single word, finds a satisfying form
for his cognition means, by his word, much more than is uttered thereby. Were he in
a position to bring the one-sidedness of his cognition to clear consciousness, he
would with the same stroke have come to the insight that he had to proceed to an
opposite determination, which had just as universal a meaning. But the mutually
opposed determinations are united in a third, more correct category, in which the
two preceding ones are both preserved in their relative validity. The third, fuller
category is, by continued thinking, likewise recognized as one-sided, passes over
into its opposite, and thereupon forms, in union with it, a new, more complete
category.

(By these few words a preliminary representation is given of the method of
ontology, which will gain clarity through the presentation of the science itself.
Besides, I intend, partly in the course of the exposition, partly at its
conclusion, to permit myself a number of elucidating reflections, with the aim of
bringing the method to clear consciousness.)

By such a movement of thought one comes to a system \opage{201}of a priori
concepts, which, when the beginning is thought as the concept of all actuality,
retains in a more developed form the universal meaning, or presents a cognition
``\textit{sub specie æternitatis}.'' In such a system of concepts an original
cognition of actuality is uttered, because therein is determined what must always
and necessarily be thought in actuality, however it may otherwise present itself.
Such an a priori circle of thought, which constitutes the content of ontology,
gives the result that personality is actuality thought in its truth, or that
nothing is with truth to be called actual except what has personality. But only in
the movement through the whole cycle of single determinations does it become clear
what actuality is and what personality is; for at this place we have only
presupposed that each of these words itself carries an immediate meaning. By such
an a priori movement of thought a twofold gain is achieved. In the first place,
the a priori concept of actuality is completely explained. Account is namely given
of how actuality must be thought always, or under the form of eternity: account is
given, in such a way, of the fact that it is only at the conclusion of the
exposition that it is uttered what one understands by the proposition that
actuality is personality. This proposition must, at the conclusion of the
exposition, have won a different meaning from the one with which it can here be
uttered without preparation. --- In the second place, the single categories are
with the same stroke explained, when they are, in the course of the exposition,
defined in their place and by their place in the whole connection. For one who,
after having thought through each fundamental concept according to its meaning in
the ontological circle of thought, turns again to empirical and to practical
thinking, these fundamental concepts, although they are used in the same
combinations in which the natural tact of language chose them, have acquired a
deeper meaning; for even to the sphere of daily life he now \opage{202}brings the
light with which ontology has pervaded them.\footnote{Ontology is the doctrine of
the eternal form of thinking and of existence: it presents, namely, the universal
determinations that hold for everything that is, which All at the same time --- as
the following presentation of the science will show --- carries with it the
determinations with which we have here and there, in the foregoing, by
anticipation designated it, namely existence, actuality. In the development and
explanation of the constantly valid determinations for the actual, God and the
world and God's relation to the world are at the same time thought in abstract
concepts, i.e.\ thought in general, without therefore the whole concrete content
of actuality being taken up into the ontological form of thought. Just as the
thinker in a single category has often thought All under the form of eternity, but
in such a way that his thinking in an abstract manner embraces therein all
actuality, so too is the whole of existence thought abstractly in the whole system
of a priori concepts, which indeed with objective validity determines all
actuality, but does not present the whole richness of existence. (Without granting
any individual arbitrary power to determine the usage that is to be used in the
free republic of philosophy, we use this many-meaninged word --- abstract --- here
in a sense that has sufficient warrant in the history of science. It is used in
science now of the residue of intuition that is retained in representation, now,
as here, of the ontological concepts in their totality, now of a single category
in which fewer determinations of thought are contained, insofar as it is compared
with a category richer in content. Aristotle designates by τὸ χωριστόν both one of
the categories thought for itself and the concept drawn off from particular
objects of sense.) --- [In the margin of the author's manuscript is noted here:
``This § must be considered more closely, especially the usage: Ontology and
Metaphysics.'']}

Ontology utters, as has been said, one principal moment of \opage{203}scientific
cognition: but the second principal moment of the latter is experience. Only the
experience that has the thoroughly worked-through system of categories for its
presupposition, or the a priori thinking that is united with sympathy and with the
cognition of actual persons, is the energy of human cognition. The one-sided a
priori thinking of thinking is not divine, but in worth below the true human
thinking. Experience without pure thinking is blind, and an a priori concept
without personally lived experience is mere formal conceiving or
schematism.\footnote{This piece may perhaps be developed further, especially the
relation between pure thinking and experience. \emph{The author's marginal note.}}

Ontology has this twofold aim: partly to bring to consciousness the more abstract
determinations that constitute the content of the more concrete categories, partly
to show how those determinations unite themselves with necessity in the more
concrete ones, so that the movement of thought deserves the name of both analytic
and synthetic; for it is precisely by the separating thinking, the separation of a
single abstract determination, that its inseparability from the opposite
determination, by which it is integrated, is most clearly seen. But here the
question arises how far the analysis is to go; for one person and another seem to
take it as a settled matter that the most perfect thinking here separates the most
moments --- naturally under the presupposition that the synthesis constantly
accompanies it. That the separation here does not strike upon any absolute
boundary is easily seen, since the finer thinking constantly has it in its power to
insert more members into the series of the discrete concepts. But since thinking is
nevertheless only genuine (true) thinking by being definitely uttered, it can, from
want of fitting words, easily come to drag itself along with a sham life. The
thinker can easily come to lose himself in hairsplitting, where the word entirely
steps into the \opage{204}place of the supposed thought. Just as speech, by
imitating the form of the movement of thought, can be continued after thinking has
ceased to move within it, so too, in the transition to empty speech, a flat and
half-dead speech can be produced. Its proper, fresh life both philosophical
thinking and other thinking has only so long as the immediate meaning of the word
in the mother tongue is not entirely abandoned. From this the Greek philosophy has
its fresh colors; for it grew forth out of the popular thinking of the people (the
tribe), and it is borne by the immediate meaning of the words as by a vigorous
root. Weaker becomes the thinking that proceeds within foreign tribes' word-forms
and comes forward in a motley dress of foreign turns of speech. Yet they may still,
for the learned man who has entered into foreign peoples' idiom, have much derived
life left in them. But when individual arbitrariness in the formation of words
gains all too great scope, speech loses ever more and more of its meaning, because
it spreads itself out in the air, torn loose from its root and its life-principle.
Longest it continues to have validity, as a merely subjective form, for the thinker
himself, but at last the solitary thinking, in self-made forms, becomes mere
word-mongering. Herein lies one of the grounds why the analyzing thinking can go
too far; it can namely approach too near the point where the word becomes partly an
arbitrary, merely subjective expression for the thought, partly wholly thoughtless.
The old proposition, \textit{qui bene distinguit, bene docet}, is therefore rather
one-sided, unless one lays particular emphasis on the addition \textit{bene}. By
this \textit{bene} it must then partly be expressed that the distinction must
constantly go hand in hand with the conjunction, partly that the analysis must not
be continued until it becomes a contentless, merely sham thinking.

But that a presentation of philosophy (ontology) has overstepped the proper
boundaries for the articulation of the concept and has \opage{205}lost itself in
subtleties can only be judged immediately by the discerning connoisseur; there
naturally lets itself be adduced no scientific proof for it. The proper boundary
here becomes, according to the demands of the individualities (the people) and of
the ages, of a highly variable nature, so that herein lies one of the many grounds
that compel philosophy to come forward in a manifoldness of individual forms. The
presentation of philosophy with which an individual, without affectation, finds
himself satisfied, and in which he moves freely, is now more concentrated, now more
developed; but a single norm valid for all can never be found. We here set up only
the general proposition, that the separating out of the abstract determinations
that are united in a concept can be driven too far. Treatises in which this fault
is committed therefore attract little attention in the literature: they have,
namely, interest only for the writer himself, but not for the reader. Such an
exaggerated analysis is not seldom undertaken by Hegelians, who, after Hegel's own
method, continue the dismemberment of the concepts at which he comes to a stop.
`How lifeless is the presentation in several of Hinrichs's works! Even Hegel, who
yet had most occasion to be pleased with it, begs, in his letters, the author to
lay greater diligence upon clarity, and confesses that it often falls rather
troublesome to him himself to follow his exposition.

Just as the presentation of philosophy often loses its intensity in that the
concept is all too much dissolved into its abstract components and, so to speak,
unraveled into thin threads, so it also often suffers from the opposite defect and
thereby loses its objective validity; we think here of the superficial exposition
that makes all too long strides in its dialectic, skips over too many intermediate
members, and thereby becomes unintelligible to the reader. But the proper measure
here too depends on the time and the definite circumstances under which the
presentation has come to be. --- Thus the definite \opage{206}relation between
intension and extension in the formation of philosophical thinking is one of the
ways in which the national spirit expresses itself among a people. --- ---

\bigskip
\begin{center}
{\large\scshape Chapter One.}\\[0.4em]
{\large\bfseries The Doctrine of the Single Concepts.}
\end{center}
\begin{center}\rule{0.25\textwidth}{0.4pt}\end{center}
\begin{center}
{\bfseries First Section.}\\[0.3em]
{\bfseries The Beginning-Concepts.}
\end{center}
\medskip

By reflections on the nature of cognition, the natural certainty of the objective
reality of thinking can be disturbed to such a degree that all cognition of truth
seems to transform itself into a meaningless play with merely subjective concepts.
Thereby arise such shapes of consciousness as idealism, skepticism, and nihilism,
which have their common ground in this, that the thinker first, by a violent
abstraction, wholly separates thinking from existence, and then in vain looks
around for grounds of their agreement.

But in continuing the train of thought that regards the fundamental concepts of
thinking as merely subjective forms, one notices that it cannot possibly be carried
through with consistency. It cannot be uttered at all except under the
presupposition of the real validity of certain fundamental concepts, of which one
can make use as means to its presentation. But once one has made this clear to
oneself, the merely subjective thinking is overcome; the self-contradiction in
which it at last entangles itself has been brought to consciousness, and the way is
paved for an advance in cognition.

The concepts by which existence is thought cannot, namely, \opage{207}possibly all
together lose their meaning for the doubter. This is experienced all the more
clearly the more one approaches, with one's skeptical reflection, the most abstract
among them. We will at once direct attention to the most contentless fundamental
concept found in our consciousness, namely the concept Being, or to be; for therein
we meet precisely an insurmountable boundary for idealistic doubts and broodings.
The opinion that the concept Being has no objective validity no one has consistently
defended, and no one can clearly utter it; for if one does not presuppose the reality
of this concept, one can assert of nothing that it is, or that it is not. From this
comes the often-emphasized contradiction in which Kant's critique of reason is
entangled. According to the plan of the system, Being had to be regarded as a
subjective concept that held only for the world of phenomena; but notwithstanding,
there is continual talk of that which, in itself, is, without the least account
being given of whence the concept of it is fetched, and whence it acquires its
reality. But it is the nature of the matter that has here made itself valid against
a contrived train of thought: one can assert nothing and deny nothing without
ascribing to the concept Being more than subjective meaning; indeed, without
assuming its reality one has no means of denying it its reality. --- This concept,
which carries with it such invincible certainty of reality, is highly fitted to make
thinking's indissoluble connection with Being quite evident, and this concept we
have, as has been said, chosen as the point of departure for ontology.

\emph{A}~~\emph{Being.} Since ontology is here expounded without its scientific
introduction, the concept Being is here taken as immediate. It is presupposed that
the word carries the meaning in which it is here to be taken; but in the advance to
the more concrete concepts, into which this beginning-concept will dialectically
develop itself, \opage{208}this meaning will come forward more definitely, in that
the transition itself here, as with all the following categories, will cast light
back upon the foregoing.

The thought Being is the beginning-point in the properly philosophical or productive
thinking, and it is the aim of science to let a whole a priori system of thought
arise through the expansion of this point. In that the concept Being develops itself,
for thinking, into concepts richer in content, thinking can with truth be said itself
to bring forth its content, and the content that it itself brings forth is a cycle of
determinations for all actuality, which have universal validity and necessity. ---
Sometimes the assertion is disputed that one can, by an a priori or, as one now often
hears, immanent thinking, develop the most universal determinations of thinking and
of existence in such a way that thinking, after abstraction from all content, with
free self-sufficiency reproduces itself in the true form that belongs to it. The
objections against this doctrine rest entirely on misunderstanding. But this
misunderstanding lets itself be removed by a couple of words, which at the same time
can contribute to drawing off some of the aristocratic obscurity from the concept of
immanent thinking. We naturally grant that the fundamental concepts of thinking
develop themselves under experience, and that the most important categories, which
constitute the content of ontology, must be found in the thinkers' consciousness
before the idea of ontology can be grasped. The production spoken of is therefore, in
a certain sense, a more thorough reproduction, by which one partly makes clear to
oneself the relation that obtains among those fundamental concepts, of which one had
indeed previously made an immediate use, partly --- and this is the most important
--- in that one orders them, finds them to be interconnected fragments of an original
human fundamental cognition. By this altered position and meaning, which the
categories acquire after that reproduction, one is thus really justified in regarding
it as a \opage{209}production of something new. --- Besides, the ontologist is by no
means restricted to the cycle of categories that are handed down to him as expressly
named in language. While he, in a methodical way, surveys the connection of the
categories, it may well be possible that the science requires a different articulation
from the one that the need of social life has occasioned. In every thorough ontology
there are therefore found categories that are not known outside the scientific
language. But it is here that the airy formations of words and formulas can so easily
gain entrance, in that the lust for mastery and subjective arbitrariness make
themselves valid in the choice of technical terms. Only when this awkward state comes
to an end will the study of philosophy become truly gladdening, because human thinking
can come to move in forms that accord with the laws of nature and of beauty. Only then
will the true priests of philosophy come to conduct its word. --- Who indeed doubts
that nowadays those who are properly called by nature to philosophical research find
themselves shut out from it by the prevailing philosophers' narrowness, want of the
gift of presentation, and dominant guild-spirit. Philosophical genius is surely not
the foremost requisite for one who is to be able to orient himself in the factual
condition of the now-prevailing German philosophy.

As the beginning of all determination of thought,\footnote{Here the author has noted
the following: ``Partly written during my illness this year (1838) and begun during
asthmatic attacks last year, after another failed draft had been discarded.''} the
concept Being must of all be the most abstract; no various determinations let
themselves be distinguished within it as its content; it is wholly simple and
indissoluble. Therefore it is impossible to give any definition of it in the usual
meaning of the word Definition. For that, namely, it would be required that there be
given a more abstract \opage{210}concept under which it could be subsumed; but then
this concept would have to be regarded as a more correct point of departure for
ontological thinking than the concept Being. The concept Being becomes defined only
by its place in the whole system of the categories. A consequence of the abstract
nature of Being is its universal validity. Everything that lets itself be determined
by any other category, such as cause or consequence, or substance, etc., also lets
itself be determined by this most universal of all concepts, which lies at the ground
of them all. When one of the more concrete categories is predicated of a subject,
Being is thereby indirectly predicated of it; whereas by mere Being not a single one
of the other determinations of thought is uttered of a subject.

Since Being in an abstract manner encloses all actuality and is contained in all other
determinations of thought for actuality, the concept is thought incorrectly if it is
apprehended as an opposite to thinking. For thus, according to its immediate meaning,
the word can also be understood. But according to the all-embracing meaning in which
Being is here to be taken, it holds of all that is subjective as well as of all that is
objective. Thinking too falls under the determination Being.

As a transition to real thinking, or beginning thereto, Being is at once both
determination and non-determination. This contradiction makes it so that one cannot
remain standing at this concept, but is driven from it to other concepts by which it is
completed. The need for closer determination that accompanies the concept Being we will
strive to bring to clearer consciousness by yet a couple of different turns of phrase.

Being can be regarded as the mere copula with a wholly indeterminate subject that seeks
its predicate. It utters neither that something is, nor what it is. \textit{x} is
\textit{x} is the formula in which this complete want of determinateness can be uttered.
By this \opage{211}reflection, that we here have, of the form of the proposition, only
the mere copula --- is ---, it becomes quite intuitable that we have to do with an
unfinished thought that demands its supplement.

The wholly indeterminate Being is as yet Nothing. Thus one is oneself, even according to
immediate usage, justified in expressing oneself. What cannot be determined by any
predicate, what is not Something, that is Nothing. (This is one of the several senses in
which the word Nothing can be taken.) By the consideration that Being is Nothing, or
rather, that it is as yet for us Nothing, the ontological movement of thought has
properly first begun. For it presupposes, like every movement, two points: one from
which the movement proceeds, and one to which it goes. By this movement it is here
brought to definite consciousness that in the concept Being there lies a lack or a want.
Nothing here designates the vacant place for a predicate. The pure Being is Nothing says
the same as: to be is as yet Nothing, unless thereby is thought to be Something.
\emph{Something} is thus the new category that steps into the place of Nothing, and only
by uttering the proposition ``\emph{Being is Something}'' have I made a clear advance in
thinking.

It has become customary, on Hegel's authority, to regard the pure Being as the category
that lies at the ground of the Eleatics' worldview. Without deciding the question whether
that philosophy is with good ground compared with this place of ontology (which, from a
strictly historical standpoint, would perhaps have to be denied), it is remarked in
passing that τὸ ὄν among the Eleatics --- instead of τὸ εἶναι --- does not, at least, quite
correspond to the pure Being or the to-be, since the hypostatization that lies in the
Greek word makes the concept more concrete than it here, according to its determination,
ought to be. The Danish \emph{Væren} (Being) is here more adequate than the corresponding
words among several \opage{212}other cultivated nations. The French \textit{être} is,
e.g., more ambiguous, as is the German \emph{Seyn}, which, at least with the definite
article, is also used to designate the sum-total of all that is.

\emph{B.}~~\emph{Something.} The proper advance from the indeterminate Being is the one
that is made to determinateness in general. To be is to be determined in one way or
another, or: All that is, is determined, is Something --- ---

\begin{center}\rule{0.32\textwidth}{0.4pt}\end{center}

\medskip
\begin{center}\footnotesize\textit{[The fragment breaks off here.]}\end{center}

\end{document}
